% GENERATED FILE. Edit canonical lessons, then run npm run generate:books.
% canonical-source-hash: fnv1a64:9300ee08abb0ded6
% canonical-lessons: KA-C456-purabhavana, KA-C456-gopura, KA-C456-karanji, KA-C456-heggurutu, KA-C456-viksisu

\chapter{Sights in the Old Town}
\label{ch:ka-sights-in-the-old-town}
\hlchaptermodality{456}

\begin{chapteropening}
\textbf{By the end of this chapter:} \emph{I can say a town hall, a tower, a fountain, a landmark and to view.}

Complete the last lesson of chapter 456: I can say a town hall, a tower, a fountain, a landmark and to view.
\end{chapteropening}

\section[\texorpdfstring{\kn{ಪುರಭವನ} (purabhavana)}{ಪುರಭವನ (purabhavana)}]{\kn{ಪುರಭವನ} (purabhavana) — a town hall}

\label{lesson:KA-C456-purabhavana}



\begin{quote}
Before the new one: say the Kannada for a guide, then the Kannada for worth seeing.
\end{quote}

\subsection*{You'll want to know: \kn{ಪುರಭವನ}}
\textbf{\kn{ಪುರಭವನ}} — \emph{purabhavana} — \textquotedblleft{}a town hall\textquotedblright{}.

\begin{quote}
\kn{ಪುರಭವನದ} \kn{ಮುಂದೆ} \kn{ಒಂದು} \kn{ದೊಡ್ಡ} \kn{ಗಡಿಯಾರ} \kn{ಇದೆ}.
\end{quote}

\begin{quote}
\emph{purabhavanada munde ondu doḍḍa gaḍiyāra ide.}
\end{quote}

\begin{quote}
There is a big clock in front of the town hall.
\end{quote}

\kn{ಪುರ}, a town, plus \kn{ಭವನ}, a building or mansion.

\begin{cousinweb}
\textbf{\kn{ಪುರಭವನ}} (\emph{purabhavana}) is a modern compound built from Sanskrit words: \emph{pura}, \textquotedblleft{}a fortress, a town\textquotedblright{} (CDIAL 8278), and \emph{bhavana}, \textquotedblleft{}a place of abode\textquotedblright{} (CDIAL 9417).
\end{cousinweb}

\begin{grammarlens}[title={What you've built}]
One more word for getting around.
\end{grammarlens}

\subsection*{Your turn}
\begin{itemize}
\raggedright
  \item \emph{Say it:} \emph{purabhavana}
  \item \emph{Say it:} \emph{purabhavana}, once more
  \item \emph{Say it:} \emph{purabhavanada munde ondu doḍḍa gaḍiyāra ide.} — There is a big clock in front of the town hall.
  \item \emph{Recall:} say the Kannada for a guide, then the Kannada for worth seeing, then say \emph{purabhavana} again
\end{itemize}

\begin{tcolorbox}[breakable,colback=teal!4,colframe=teal!35!black,title={Before you move on}]
What does \kn{ಪುರಭವನ} mean? (\textquotedblleft{}A town hall\textquotedblright{}.) Say it once more.
\end{tcolorbox}
\section[\texorpdfstring{\kn{ಗೋಪುರ} (gōpura)}{ಗೋಪುರ (gōpura)}]{\kn{ಗೋಪುರ} (gōpura) — a tower; a temple tower}

\label{lesson:KA-C456-gopura}



\begin{quote}
Before the new one: say the Kannada for worth seeing, then the Kannada for a town hall.
\end{quote}

\subsection*{You'll want to know: \kn{ಗೋಪುರ}}
\textbf{\kn{ಗೋಪುರ}} — \emph{gōpura} — \textquotedblleft{}a tower; a temple tower\textquotedblright{}.

\begin{quote}
\kn{ದೇವಸ್ಥಾನದ} \kn{ಗೋಪುರ} \kn{ದೂರದಿಂದಲೇ} \kn{ಕಾಣುತ್ತದೆ}.
\end{quote}

\begin{quote}
\emph{dēvasthānada gōpura dūradindalē kāṇuttade.}
\end{quote}

\begin{quote}
The temple tower can be seen from far away.
\end{quote}

\kn{ಗಡಿಯಾರ} \kn{ಗೋಪುರ}, a clock tower.

\begin{cousinweb}
\textbf{\kn{ಗೋಪುರ}} (\emph{gōpura}) is borrowed from Sanskrit \emph{gōpura}, \textquotedblleft{}a town gate\textquotedblright{} (Monier-Williams); in South India it names the tall tower over a temple gate.
\end{cousinweb}

\begin{grammarlens}[title={What you've built}]
One more word for getting around.
\end{grammarlens}

\subsection*{Your turn}
\begin{itemize}
\raggedright
  \item \emph{Say it:} \emph{gōpura}
  \item \emph{Say it:} \emph{gōpura}, once more
  \item \emph{Say it:} \emph{dēvasthānada gōpura dūradindalē kāṇuttade.} — The temple tower can be seen from far away.
  \item \emph{Recall:} say the Kannada for worth seeing, then the Kannada for a town hall, then say \emph{gōpura} again
\end{itemize}

\begin{tcolorbox}[breakable,colback=teal!4,colframe=teal!35!black,title={Before you move on}]
What does \kn{ಗೋಪುರ} mean? (\textquotedblleft{}A tower; a temple tower\textquotedblright{}.) Say it once more.
\end{tcolorbox}
\section[\texorpdfstring{\kn{ಕಾರಂಜಿ} (kārañji)}{ಕಾರಂಜಿ (kārañji)}]{\kn{ಕಾರಂಜಿ} (kārañji) — a fountain}

\label{lesson:KA-C456-karanji}



\begin{quote}
Before the new one: say the Kannada for a town hall, then the Kannada for a tower.
\end{quote}

\subsection*{You'll want to know: \kn{ಕಾರಂಜಿ}}
\textbf{\kn{ಕಾರಂಜಿ}} — \emph{kārañji} — \textquotedblleft{}a fountain\textquotedblright{}.

\begin{quote}
\kn{ಉದ್ಯಾನದ} \kn{ನಡುವೆ} \kn{ಒಂದು} \kn{ಸುಂದರ} \kn{ಕಾರಂಜಿ} \kn{ಇದೆ}.
\end{quote}

\begin{quote}
\emph{udyānada naḍuve ondu sundara kārañji ide.}
\end{quote}

\begin{quote}
There is a beautiful fountain in the middle of the park.
\end{quote}

Many town circles have a \kn{ಕಾರಂಜಿ}.

\begin{cousinweb}
\textbf{\kn{ಕಾರಂಜಿ}} (\emph{kārañji}) is a Kannada word whose origin is not settled.
\end{cousinweb}

\begin{grammarlens}[title={What you've built}]
One more word for getting around.
\end{grammarlens}

\subsection*{Your turn}
\begin{itemize}
\raggedright
  \item \emph{Say it:} \emph{kārañji}
  \item \emph{Say it:} \emph{kārañji}, once more
  \item \emph{Say it:} \emph{udyānada naḍuve ondu sundara kārañji ide.} — There is a beautiful fountain in the middle of the park.
  \item \emph{Recall:} say the Kannada for a town hall, then the Kannada for a tower, then say \emph{kārañji} again
\end{itemize}

\begin{tcolorbox}[breakable,colback=teal!4,colframe=teal!35!black,title={Before you move on}]
What does \kn{ಕಾರಂಜಿ} mean? (\textquotedblleft{}A fountain\textquotedblright{}.) Say it once more.
\end{tcolorbox}
\section[\texorpdfstring{\kn{ಹೆಗ್ಗುರುತು} (heggurutu)}{ಹೆಗ್ಗುರುತು (heggurutu)}]{\kn{ಹೆಗ್ಗುರುತು} (heggurutu) — a landmark}

\label{lesson:KA-C456-heggurutu}



\begin{quote}
Before the new one: say the Kannada for a tower, then the Kannada for a fountain.
\end{quote}

\subsection*{You'll want to know: \kn{ಹೆಗ್ಗುರುತು}}
\textbf{\kn{ಹೆಗ್ಗುರುತು}} — \emph{heggurutu} — \textquotedblleft{}a landmark\textquotedblright{}.

\begin{quote}
\kn{ಆ} \kn{ಗೋಪುರ} \kn{ನಗರದ} \kn{ಹೆಗ್ಗುರುತು}.
\end{quote}

\begin{quote}
\emph{ā gōpura nagarada heggurutu.}
\end{quote}

\begin{quote}
That tower is the landmark of the city.
\end{quote}

\kn{ಹೆ}-, big, plus \kn{ಗುರುತು}, a mark: the big sign people steer by.

\begin{cousinweb}
\textbf{\kn{ಹೆಗ್ಗುರುತು}} (\emph{heggurutu}) is a native Kannada compound of \emph{per}, \emph{hē-}, \textquotedblleft{}big\textquotedblright{} (DEDR 4411), and \kn{ಗುರುತು}, \textquotedblleft{}a mark, a sign\textquotedblright{} (DEDR 1847).
\end{cousinweb}

\begin{grammarlens}[title={What you've built}]
One more word for getting around.
\end{grammarlens}

\subsection*{Your turn}
\begin{itemize}
\raggedright
  \item \emph{Say it:} \emph{heggurutu}
  \item \emph{Say it:} \emph{heggurutu}, once more
  \item \emph{Say it:} \emph{ā gōpura nagarada heggurutu.} — That tower is the landmark of the city.
  \item \emph{Recall:} say the Kannada for a tower, then the Kannada for a fountain, then say \emph{heggurutu} again
\end{itemize}

\begin{tcolorbox}[breakable,colback=teal!4,colframe=teal!35!black,title={Before you move on}]
What does \kn{ಹೆಗ್ಗುರುತು} mean? (\textquotedblleft{}A landmark\textquotedblright{}.) Say it once more.
\end{tcolorbox}
\section[\texorpdfstring{\kn{ವೀಕ್ಷಿಸು} (vīkṣisu)}{ವೀಕ್ಷಿಸು (vīkṣisu)}]{\kn{ವೀಕ್ಷಿಸು} (vīkṣisu) — to view; to watch}

\label{lesson:KA-C456-viksisu}



\begin{quote}
Before the new one: say the Kannada for a fountain, then the Kannada for a landmark.
\end{quote}

\subsection*{You'll want to know: \kn{ವೀಕ್ಷಿಸು}}
\textbf{\kn{ವೀಕ್ಷಿಸು}} — \emph{vīkṣisu} — \textquotedblleft{}to view; to watch\textquotedblright{}.

\begin{quote}
\kn{ಪ್ರವಾಸಿಗರು} \kn{ಬೆಟ್ಟದಿಂದ} \kn{ಸೂರ್ಯಾಸ್ತವನ್ನು} \kn{ವೀಕ್ಷಿಸಿದರು}.
\end{quote}

\begin{quote}
\emph{pravāsigaru beṭṭadinda sūryāstavannu vīkṣisidaru.}
\end{quote}

\begin{quote}
The tourists watched the sunset from the hill.
\end{quote}

Past (I …): \kn{ವೀಕ್ಷಿಸಿದೆ}, \emph{vīkṣiside}.

A formal word for \kn{ನೋಡು}: \kn{ಟೀವಿ} \kn{ವೀಕ್ಷಕರು}, TV viewers.

\begin{cousinweb}
\textbf{\kn{ವೀಕ್ಷಿಸು}} (\emph{vīkṣisu}) is built from Sanskrit \emph{vīkṣ}, \textquotedblleft{}to behold, to see\textquotedblright{} (CDIAL 12041, \emph{vīkṣatē}), with the Kannada verb ending -\emph{isu}.
\end{cousinweb}

\begin{grammarlens}[title={What you've built}]
One more word for getting around.
\end{grammarlens}

\subsection*{Your turn}
\begin{itemize}
\raggedright
  \item \emph{Say it:} \emph{vīkṣisu}
  \item \emph{Say it:} \emph{vīkṣisu}, once more
  \item \emph{Say it:} \emph{pravāsigaru beṭṭadinda sūryāstavannu vīkṣisidaru.} — The tourists watched the sunset from the hill.
  \item \emph{Recall:} say the Kannada for a fountain, then the Kannada for a landmark, then say \emph{vīkṣisu} again
\end{itemize}

\begin{tcolorbox}[breakable,colback=teal!4,colframe=teal!35!black,title={Before you move on}]
What does \kn{ವೀಕ್ಷಿಸು} mean? (\textquotedblleft{}To view; to watch\textquotedblright{}.) Say it once more.
\end{tcolorbox}
